\chapter{Extremalelemente und Auswahlwerte}
\label{b15guide:ch:extrema}

\section{Minimale und kleinste Elemente}

\begin{definition}[Minimalität und Maximalität]
Für $S\subseteq A$ und $m,M\in S$ heißt $m$ \emph{minimal
in $S$}, wenn kein $s\in S$ die Bedingung $s<m$ erfüllt.
$M$ heißt \emph{maximal in $S$}, wenn kein $s\in S$
die Bedingung $M<s$ erfüllt.
\end{definition}

\begin{theorem}[Minimal ist kleinstes, maximal ist größtes]
\label{b15guide:thm:minimal}
Für $S\subseteq A$ gilt:
\[
\begin{aligned}
 m\text{ ist minimal in }S&\Longleftrightarrow m=\min S,\\
 M\text{ ist maximal in }S&\Longleftrightarrow M=\max S.
\end{aligned}
\]
Auf der rechten Seite ist die Existenz des jeweiligen Wertes
eingeschlossen. Insbesondere hat jede Teilmenge von $A$
höchstens ein minimales und höchstens ein maximales Element.
\end{theorem}

\begin{theorem}[Fehlen eines Endpunkts]
\label{b15guide:thm:no-endpoint}
Für eine nichtleere Teilmenge $S\subseteq A$ gilt
\[
\begin{aligned}
 S\text{ hat kein Minimum}
 &\Longleftrightarrow
   \forall x\in S\,\exists y\in S\;(y<x),\\
 S\text{ hat kein Maximum}
 &\Longleftrightarrow
   \forall x\in S\,\exists y\in S\;(x<y).
\end{aligned}
\]
Die Aussage fordert jeweils ein passendes Element;
eine gleichzeitig gewählte Folge solcher Elemente wird nicht behauptet.
\end{theorem}

\begin{theorem}[Endliche Teilmengen sind eindeutig sortierbar]
\label{b15guide:thm:sorting}
Jede nichtleere endliche Teilmenge $S\subseteq A$ mit $n$ Elementen
besitzt eine eindeutig bestimmte Aufzählung
\[
 S=\{s_1,\ldots,s_n\},\qquad s_1<s_2<\cdots<s_n.
\]
Dabei gilt $\min S=s_1$ und $\max S=s_n$.
Für $n=1$ fallen beide Werte zusammen.
\end{theorem}

\section{Minimum und Maximum zweier Elemente}

\begin{definition}[Binäre Minimums- und Maximumswerte]
\label{b15guide:def:minmax}
Für $x,y\in A$ setzen wir
\[
 \min(x,y)=\min\{x,y\},\qquad
 \max(x,y)=\max\{x,y\}.
\]
Die Schreibweisen kürzen
$\min_{A,\leq}(x,y)$ und $\max_{A,\leq}(x,y)$ aus Band~15 ab.
Bei endlichen nichtleeren Mengen wird entsprechend $\min S$
beziehungsweise $\max S$ geschrieben.
\end{definition}

\begin{theorem}[Existenz und Auswahl eines Operanden]
\label{b15guide:thm:selection}
Für jedes Paar $x,y\in A$ existieren beide Werte, und es gilt
\[
\begin{aligned}
 \min(x,y)&=
 \begin{cases}x,&x\leq y,\\ y,&y<x,\end{cases}
&
 \max(x,y)&=
 \begin{cases}y,&x\leq y,\\ x,&y<x.\end{cases}
\end{aligned}
\]
Insbesondere liegen beide Werte in $\{x,y\}$, und
\[
 \{\min(x,y),\max(x,y)\}=\{x,y\}.
\]
\end{theorem}

\begin{theorem}[Vergleichsregeln mit einer Disjunktion]
\label{b15guide:thm:or-comparisons}
Für $x,y,z\in A$ gelten
\[
\begin{aligned}
 \min(x,y)\leq z&\Longleftrightarrow(x\leq z\lor y\leq z),\\
 \min(x,y)<z&\Longleftrightarrow(x<z\lor y<z),\\
 z\leq\max(x,y)&\Longleftrightarrow(z\leq x\lor z\leq y),\\
 z<\max(x,y)&\Longleftrightarrow(z<x\lor z<y).
\end{aligned}
\]
Die entsprechenden universellen Vergleichsregeln mit
Konjunktionen gehören bereits zum allgemeinen Paarwertkalkül
in Band~14 und werden hier nicht erneut aufgeführt.
\end{theorem}

\begin{theorem}[Endliche Auswahl und verschachtelte Werte]
\label{b15guide:thm:iterated}
Für $n\geq1$ und $x_1,\ldots,x_n\in A$ liefert jeder vollständig
geklammerte Ausdruck aus binären Minima, in dem jedes der
aufgelisteten Elemente einmal vorkommt, den Wert
\[
 \min\{x_1,\ldots,x_n\}\in\{x_1,\ldots,x_n\}.
\]
Entsprechend liefern solche Ausdrücke aus Maxima den größten
der aufgelisteten Werte. Die Reihenfolge und Klammerung ändern
den gewählten Wert nicht; Wiederholungen von Werten sind erlaubt.
Für $n=1$ ist der Ausdruck das einzelne Element.
\end{theorem}

\begin{theorem}[Endliche disjunktive Vergleichsregeln]
\label{b15guide:thm:finite-or}
Für nichtleeres endliches $S\subseteq A$ und $z\in A$ gilt
\[
\begin{aligned}
 \min S\leq z&\Longleftrightarrow\exists s\in S\;(s\leq z),&
 \min S<z&\Longleftrightarrow\exists s\in S\;(s<z),\\
 z\leq\max S&\Longleftrightarrow\exists s\in S\;(z\leq s),&
 z<\max S&\Longleftrightarrow\exists s\in S\;(z<s).
\end{aligned}
\]
\end{theorem}

